% !TITLE 相见欢·无言独上西楼
% !AUTHOR 李煜
% !DYNASTY 五代十国
% !GENRE 词
% !ORDER 6
% !EDITION 《南唐二主词笺注》，王仲闻校订、陈书良与刘娟笺注，中华书局（2013年版为候选）
% !EDITION_STATUS pending

\begin{poem}
无言独上西楼，月如钩\textsuperscript{1}。
寂寞梧桐深院锁清秋\textsuperscript{2}。

剪不断\textsuperscript{3}，理还乱\textsuperscript{4}，是离愁\textsuperscript{5}。
别是一般\textsuperscript{6}滋味在心头。
\end{poem}

\begin{annotations}
\notetext{1}{月如钩：弯月像一把钩子。残月如钩，本身就带着残缺和锐利的感觉，钩住了愁绪。}
\notetext{2}{锁清秋：清冷的秋意被深锁在院子里。一个“锁”字，把无形的秋意变成了有形的囚笼——被锁住的不只是秋天，更是词人自己。}
\notetext{3}{剪不断：想剪断它，却怎么也剪不断。把离愁比作一根无形的绳索。}
\notetext{4}{理还乱：想整理它，反而更乱。和上句呼应——愁绪越想摆脱越缠绕。}
\notetext{5}{离愁：离别之愁。这里不是普通的离别，而是与故国、与从前一切美好事物的永别。}
\notetext{6}{别是一般：另有一种、不同于寻常的。这种滋味无法用语言形容——不是酸，不是苦，而是一种从未体验过的复杂感受。}
\end{annotations}

\begin{appreciation}
这首《相见欢》是李煜词里最短的几首之一，通篇不过三十六字，却把“愁”这种东西写到了底。

“无言独上西楼，月如钩。”六个字里塞进三桩事：没有话，一个人，登上西楼。无言是无人可说独是无人作伴，西楼是往高处、冷处去。“月如钩”更狠，一弯缺月尖而冷，像一枚钩子，把那点愁也一并钩住了。

“寂寞梧桐深院锁清秋。”这一句全落在一个“锁”字上。秋色本来漫天都是这里却锁进了深深的院子，院子锁住了秋，其实也锁住了院中的这个人。季节锁在墙内，人也困在墙内，一个“锁”字把无形的囚禁写得伸手可摸。

下片忽然改用短句。“剪不断，理还乱，是离愁。”三字一顿一顿一叹读起来像一声接一声的喘息。他把愁比作一团乱麻：要剪，剪不断；要理，越理越乱。这团乱麻缠的正是他与故国与从前那个自己的一场永诀。“别是一般滋味在心头”，另有一种滋味压在心底，说不出是酸是苦，只觉得那是一种自己从未尝过的味道。

我们前面读过的愁多半都是有来由的。《关雎》里“求之不得，辗转反侧”，是心上人求不到；《蒹葭》里“溯洄从之，道阻且长”，是追着跑却总差一步。这些愁都是朝前看的，前面还亮着一线指望。李煜的愁却是朝后看的，回头望去，只剩一座空落落的旧院子。他词中的月亮，在《虞美人》里照着故国，在这里弯成了一枚钩子，钩着的全是再也回不去的东西。

通篇全是寻常字眼，无一生僻，可拼到一起就成了词史上写得最透的一段愁。他写这首词的时候身在何处，那座西楼究竟在南唐还是汴京，今已难考；哪一座都好，反正他登上去，看见的都是同一钩残月。
\end{appreciation}
